\chapter{Protocolo propuesto para ensayo de laboratorio}
\label{app:ensayo_laboratorio}

Este apéndice propone un protocolo de ensayo para la etapa posterior a la simulación. Su objetivo es transformar las conclusiones del modelo en una prueba física controlada. El ensayo no pretende validar una red de campus completa; busca medir si un enlace corto puede preparar, transmitir y detectar eventos compatibles con BB84 \textit{time-bin}, y si las métricas principales se comportan como predice la simulación.

\section{Objetivo del ensayo}

El objetivo mínimo es establecer un enlace Alice--Bob de corta distancia con fibra controlada, registrar detecciones temporales, estimar QBER en una secuencia de prueba y medir parámetros físicos que luego puedan alimentar el simulador. La prueba debe priorizar trazabilidad sobre tasa máxima. Es preferible obtener pocos datos confiables, con configuración documentada, que intentar una tasa alta sin poder explicar pérdidas o errores.

El ensayo debería responder cinco preguntas:

\begin{enumerate}
  \item ¿La fuente puede operar con intensidades compatibles con una fuente coherente débil?
  \item ¿El receptor distingue bins temporales con suficiente resolución?
  \item ¿El interferómetro alcanza una visibilidad estable durante la medición?
  \item ¿El QBER medido se mantiene por debajo del umbral de referencia en una condición corta?
  \item ¿Los parámetros medidos permiten recalibrar la simulación?
\end{enumerate}

\section{Instrumental mínimo}

El instrumental depende de disponibilidad, pero la prueba debería documentar al menos fuente óptica, atenuación, fibra, receptor, detección y temporización. Si no se dispone de todos los elementos para una implementación BB84 completa, se puede realizar un ensayo parcial: primero caracterización de fuente y receptor, luego transmisión temporal, luego interferencia, y finalmente secuencia protocolar.

\begin{longtable}{p{0.24\textwidth}p{0.34\textwidth}p{0.30\textwidth}}
  \caption{Instrumental mínimo sugerido para un ensayo QKD \textit{time-bin}.}
  \label{tab:instrumental_ensayo}\\
  \toprule
  Bloque & Elementos & Medición asociada \\
  \midrule
  \endfirsthead
  \toprule
  Bloque & Elementos & Medición asociada \\
  \midrule
  \endhead
  Fuente & Láser, driver, atenuador, modulador & Potencia, estabilidad, \(\mu\) estimado \\
  Temporización & Generador de pulsos, reloj, FPGA o PC & Separación de bins, jitter \\
  Canal & Patch cords, bobinas de fibra, conectores & Pérdida total, reflexión, longitud \\
  Receptor & Switch/divisor, MZI, detectores & Visibilidad, eficiencia, dark counts \\
  Registro & Time tagger u osciloscopio rápido & Tiempos de llegada, histogramas \\
  Control & Software de adquisición y planillas & Trazabilidad, semillas, parámetros \\
  \bottomrule
\end{longtable}

\section{Procedimiento por etapas}

La primera etapa es caracterizar la fuente. Se mide potencia media, estabilidad y capacidad de atenuación. Si se busca estimar \(\mu\), debe conocerse la energía por pulso y la energía de un fotón a la longitud de onda usada. Para una longitud de onda \(\lambda\), la energía por fotón es

\begin{equation}
  E_f = \frac{hc}{\lambda}.
\end{equation}

Con potencia media \(P\) y frecuencia de repetición \(f_\text{rep}\), una estimación idealizada del número medio de fotones por pulso es

\begin{equation}
  \mu \approx \frac{P}{f_\text{rep}E_f},
\end{equation}

después de descontar atenuaciones y pérdidas conocidas. Esta estimación debe tratarse con cuidado porque depende de calibración de potencia, pérdidas internas y forma temporal de los pulsos.

La segunda etapa es medir el canal. Se registra longitud de fibra, conectores y pérdida total. Si hay bobinas de fibra, se puede comenzar con distancias cortas y agregar pérdida de forma controlada. Esta etapa se relaciona directamente con el Experimento 1 de la simulación.

La tercera etapa es caracterizar el receptor. Se mide dark count rate con la fuente apagada, respuesta temporal con pulsos conocidos y tasa de clicks con distintas intensidades. Si hay interferómetro, se mide visibilidad variando fase o condición de interferencia. Esta etapa se relaciona con los Experimentos 2 y 3.

La cuarta etapa es ejecutar una secuencia BB84 simplificada. Alice prepara una secuencia conocida o pseudoaleatoria de bits y bases. Bob registra resultados y bases. Luego se realiza tamizado y se estima QBER. Para una primera prueba, la aleatoriedad puede ser generada por software, siempre que se declare como limitación. En una versión madura, debería incorporarse una fuente aleatoria adecuada.

\section{Registro de datos}

Cada corrida de laboratorio debe tener una ficha. Esa ficha debe incluir fecha, responsables, configuración óptica, longitud de fibra, pérdidas, potencia, frecuencia, separación temporal, estado del interferómetro, detectores usados, temperatura aproximada, archivos generados y observaciones. Sin esta ficha, una medición puede ser útil para aprendizaje, pero no para sostener una memoria técnica.

\begin{table}[H]
  \centering
  \begin{tabularx}{\textwidth}{p{0.25\textwidth}X}
    \toprule
    Campo & Descripción \\
    \midrule
    ID de corrida & Código único para relacionar planilla, logs y figuras \\
    Configuración Alice & Fuente, atenuación, frecuencia, moduladores y patrón \\
    Configuración canal & Fibra, conectores, pérdida medida y longitud \\
    Configuración Bob & Detector, ventana temporal, interferómetro y visibilidad \\
    Métricas & Clicks, QBER, tasa tamizada, dark counts, observaciones \\
    Archivos & Nombre de logs, scripts, planillas y figuras asociadas \\
    \bottomrule
  \end{tabularx}
  \caption{Ficha mínima de corrida experimental.}
  \label{tab:ficha_corrida}
\end{table}

\section{Comparación con la simulación}

Después del ensayo, los parámetros medidos deberían reemplazar los valores base del simulador. La comparación no debe buscar coincidencia perfecta. Debe buscar consistencia de orden de magnitud y explicación de diferencias. Si el QBER real es mayor que el simulado, se revisan visibilidad, sincronización, dark counts, pérdidas no modeladas y asignación de bins. Si la tasa real es menor, se revisan eficiencia, \(\mu\), pérdida total y dead time.

La comparación puede organizarse en una tabla:

\begin{table}[H]
  \centering
  \begin{tabularx}{\textwidth}{p{0.24\textwidth}XXX}
    \toprule
    Magnitud & Simulación & Ensayo & Interpretación \\
    \midrule
    Pérdida total & Valor calculado & Valor medido & Ajusta transmitancia \\
    Dark counts & Parámetro fijo & Medición fuente apagada & Ajusta ruido detector \\
    Visibilidad & Valor de barrido & Medición MZI & Ajusta error de fase \\
    QBER & Resultado BB84 & Resultado tamizado & Evalúa validez del modelo \\
    Tasa & Throughput simulado & Clicks/llave medidos & Evalúa eficiencia real \\
    \bottomrule
  \end{tabularx}
  \caption{Comparación esperada entre simulación y ensayo.}
  \label{tab:sim_vs_ensayo}
\end{table}

\section{Criterios de cierre del ensayo}

El ensayo se considera exitoso si deja una línea base documentada, aunque la tasa no sea alta. Los criterios mínimos son: detección temporal repetible, estimación de QBER, medición de dark counts, medición o estimación de visibilidad, pérdida total documentada y archivos trazables. Si además el QBER queda bajo el umbral de referencia con margen, el banco está listo para una etapa de distancia mayor o para incorporar estados señuelo.

Si el ensayo falla, también produce información. Un QBER alto puede revelar problemas de alineación, temporización o interferencia. Una tasa demasiado baja puede indicar pérdidas no consideradas o mala calibración de \(\mu\). En ambos casos, el protocolo de registro permite volver al simulador con hipótesis concretas, en lugar de tratar la falla como un resultado opaco.
